
% ===========================================================================
% 4. OUTILS ET TECHNOLOGIES
% ===========================================================================
\section{Outils et technologies}

Cette section d\'ecrit les outils et technologies utilis\'es dans le projet,
ainsi que les raisons de leur choix. Les logos illustrent les technologies
officielles employ\'ees.

\subsection{Conteneurisation et d\'eploiement}

\techblock{docker}{Docker}{Docker est une plateforme de conteneurisation qui
permet d'embarquer une application et ses d\'ependances dans des conteneurs
isol\'es et reproductibles. Chaque conteneur dispose de son propre espace
processus, r\'eseau et syst\`eme de fichiers.}{La plateforme est compos\'ee de
plus de dix-huit services interd\'ependants (bases de donn\'ees, courtiers,
applications, agents de collecte). Docker garantit un d\'eploiement identique
sur toute machine, simplifie la gestion des versions et acc\'el\`ere les
red\'eploiements.}

\techblock{docker}{Docker Compose}{Docker Compose permet de d\'efinir et
d'orchestrer un ensemble de conteneurs multi-services \`a partir d'un unique
fichier YAML (r\'eseaux, volumes, limites m\'emoire, variables
d'environnement, ordre de d\'emarrage).}{Le fichier
\texttt{docker-compose.yml} centralise la d\'efinition des 18 services
applicatifs plus le serveur Fleet. C'est le point d'entr\'ee unique pour
d\'eployer et administrer l'ensemble de la plateforme.}

\subsection{D\'eveloppement}

\techblock{python}{Python 3.11}{Python est un langage de programmation
g\'en\'eraliste, interpr\'et\'e et riche d'un vaste \'ecosyst\`eme de
biblioth\`eques (parsing, HTTP, tests, calcul).}{Les modules du pipeline
(scan, enrichissement, score, indexation, tickets) sont \'ecrits en Python :
il offre une syntaxe lisible, une int\'egration native avec
\texttt{python-nmap}, \texttt{requests} et \texttt{pytest}, et acc\'el\`ere
fortement le d\'eveloppement des clients d'API (MISP, GLPI, NVD, etc.).}

\techblock{fastapi}{FastAPI}{FastAPI est un framework Python moderne et
performant pour la construction d'API REST, bas\'e sur Pydantic pour la
validation des donn\'ees et g\'en\'erant automatiquement une documentation
OpenAPI.}{Le moteur de risque expose une API \texttt{POST /score} ainsi
qu'une sonde de sant\'e \texttt{GET /health}. FastAPI a \'et\'e choisi pour sa
vitesse d'ex\'ecution, la validation automatique des entr\'ees, la
documentation interactive et sa faible latence.}

\subsection{Ordonnancement et files de t\^aches}

\techblock{celery}{Celery et Celery Beat}{Celery est un syst\`eme
d'ex\'ecution de t\^aches asynchrones distribu\'ees ; Celery Beat est son
planificateur p\'eriodique (crontab).}{Le pipeline est organis\'e en
t\^aches asynchrones (\texttt{scan}, \texttt{enrich}, \texttt{enhance},
\texttt{score}, \texttt{index\_to\_elk}, \texttt{create\_ticket})
ex\'ecut\'ees par des workers parall\`eles. Celery g\`ere la r\'epartition,
les nouvelles tentatives et le suivi des r\'esultats, tandis que Beat
d\'eclenche le scan au d\'emarrage et toutes les 6 heures, ainsi que la
synchronisation GLPI chaque heure.}

\techblock{rabbitmq}{RabbitMQ}{RabbitMQ est un courtier de messages
impl\'ementant le protocole AMQP (Advanced Message Queuing Protocol).}{Il sert
de brique centralisatrice du pipeline : chaque famille de t\^aches poss\`ede
sa file d\'ediqu\'ee (\texttt{scan}, \texttt{enrich}, \texttt{score},
\texttt{index}, \texttt{ticket}), ce qui permet de parall\'eliser le
traitement et d'absorber les pics de charge g\'en\'er\'es par les scans.}

\techblock{redis}{Redis}{Redis est une base de donn\'ees cl\'e-valeur en
m\'emoire, tr\`es rapide et compatible avec la plupart des langages.}{Redis
joue deux r\^oles : conserver l'\'etat diff\'erentiel des scans
(\texttt{voc:scan:state}) et servir de backend de r\'esultats Celery. Sa
rapidit\'e est essentielle pour comparer \`a chaque cycle des milliers de
vuln\'erabilit\'es sans p\'enaliser le pipeline.}

\subsection{Analyse de vuln\'erabilit\'es}

\techblock{nmap}{nmap}{nmap est l'outil de r\'ef\'erence pour la d\'ecouverte
de r\'eseau, l'audit de ports et la reconnaissance de services et de
syst\`emes d'exploitation.}{Il est utilis\'e avec les options
\texttt{-sS -sV -O -p- -T4} pour scanner l'int\'egralit\'e des 65535 ports
TCP, identifier les versions des services et reconna\^itre le syst\`eme
d'exploitation des h\^otes. L'acc\`es aux biblioth\`eques Python
(\texttt{python-nmap}) permet de l'int\'egrer nativement dans les t\^aches
Celery.}

\subsection{Intelligence sur les menaces}

\techblock{misp}{MISP}{MISP (\emph{Malware Information Sharing Platform}) est
une plateforme de partage d'intelligence sur les menaces : \'ev\'enements,
attributs (IOC), taxonomies et corr\'elations.}{La plateforme y recherche des
indicateurs de compromission li\'es \`a chaque CVE et y publie, \`a
l'inverse, les nouvelles vuln\'erabilit\'es sous forme d'\'ev\'enements. C'est
le pivot du volet \emph{renseignement} du pipeline.}

\techblock{nvd}{NVD}{La \emph{National Vulnerability Database} est la base de
r\'ef\'erence des CVE, g\'er\'ee par le NIST.}{L'API NVD fournit les
m\'eta-donn\'ees de chaque CVE (score CVSS, CWE, r\'ef\'erences, versions
impact\'ees) qui servent de socle \`a la corr\'elation produit/version
$\rightarrow$ CVE.}

\techblock{epss}{EPSS}{L'\emph{Exploit Prediction Scoring System} (FIRST)
estime, avec un score entre 0 et 1, la probabilit\'e qu'une vuln\'erabilit\'e
soit exploit\'ee dans les 30 jours.}{Le score EPSS est une entr\'ee du moteur
de risque : il permet de prioriser les failles r\'eellement exploitables
plut\^ot que de traiter uniquement par note CVSS.}

\techblock{kev}{CISA KEV}{Le catalogue KEV (\emph{Known Exploited
Vulnerabilities}) de la CISA recense les vuln\'erabilit\'es activement
exploit\'ees par des acteurs malveillants.}{Toute CVE pr\'esente dans ce
catalogue re\c coit un bonus de risque (2.0 par d\'efaut), car une
exploitation en cours repr\'esente une menace imm\'ediate pour le r\'eseau.}

\techblock{terminal}{OSV.dev \& Exploit-DB}{OSV.dev agr\`ege des
vuln\'erabilit\'es open-source ; Exploit-DB indexe les exploits publics
disponibles.}{Ces deux sources compl\`etent l'enrichissement : confirmation de
l'existence d'un exploit public (\texttt{exploit\_available}) et donn\'ees
suppl\'ementaires pour les composants open-source.}

\subsection{Stockage et visualisation}

\techblock{elasticsearch}{Elasticsearch}{Elasticsearch est un moteur de
recherche et d'analyse distribu\'e bas\'e sur Lucene, avec indexation
journali\`ere et agr\'egations en temps r\'eel.}{Il stocke tous les documents
de vuln\'erabilit\'es (indices \texttt{vulnerabilities-*}), les donn\'ees
d'observabilit\'e des Beats et l'\'etat des tickets (\texttt{glpi-*}). Sa
rapidit\'e et ses agr\'egations permettent des tableaux de bord
synth\'etiques.}

\techblock{logstash}{Logstash}{Logstash est un pipeline de traitement de
donn\'ees qui collecte, transforme et achemine les \'ev\'enements vers
Elasticsearch.}{Il re\c coit les documents en HTTP (\texttt{:5044}) et les
achemine vers les indices journaliers. Il centralise ainsi l'ingestion et
facilite les transformations futures sans modifier les workers.}

\techblock{kibana}{Kibana}{Kibana est l'interface de visualisation de la
stack Elastic : tableaux de bord, vues de donn\'ees, d\'ecouvertes et
administration de Fleet.}{Il fournit les tableaux de bord de
vuln\'erabilit\'es, les vues de donn\'ees (Filebeat, Packetbeat, Metricbeat,
Heartbeat, GLPI) et l'administration de Fleet. L'image a \'et\'e
personnalis\'ee pour injecter les cl\'es de chiffrement \`a
l'initialisation.}

\subsection{Gestion des agents}

\techblock{fleet}{Fleet Server et Elastic Agent}{Fleet Server est le point
d'enr\^olement central des Elastic Agents ; il distribue les politiques et
les int\'egrations de collecte.}{Fleet permet d'enr\^oler des agents sur des
h\^otes externes (Linux/Windows) via une politique par type d'h\^ote et de
superviser leur \'etat (\texttt{online}/\texttt{degraded}). C'est la brique
de \emph{gouvernance des agents}.}

\techblock{beats}{Les Beats (Filebeat, Metricbeat, Packetbeat, Auditbeat,
Heartbeat)}{Les Beats sont des agents l\'egers de collecte de donn\'ees pour
la stack Elastic.}{Chaque Beat couvre un besoin d'observabilit\'e : journaux
(Filebeat), m\'etriques syst\`eme et Docker (Metricbeat), flux r\'eseau
(Packetbeat), traces d'audit (Auditbeat) et disponibilit\'e des services
(Heartbeat). Ils alimentent les tableaux de bord de supervision.}

\subsection{Gestion des services informatiques}

\techblock{glpi}{GLPI}{GLPI est une solution ITSM open-source de gestion des
parcs informatiques, des incidents et des demandes, dot\'ee d'une API REST.}
{C'est le circuit de rem\'ediation : le pipeline y cr\'ee automatiquement un
ticket pour chaque vuln\'erabilit\'e dont le score de risque est $\geq 7$,
avec une priorit\'e et une liste de v\'erifications (\emph{checklist})
adapt\'ees \`a la s\'ev\'erit\'e.}

\subsection{Infrastructure et versionnement}

\techblock{ubuntu}{Ubuntu Server (VMware)}{Ubuntu Server est le syst\`eme
d'exploitation h\^ote de la plateforme, h\'eberg\'e sur une machine virtuelle
VMware (\texttt{192.168.184.135}).}{Il offre un \'ecosyst\`eme stable, des
paquets r\'ecents pour Docker et les outils r\'eseau (dont nmap), ainsi
qu'une documentation abondante.}

\techblock{git}{Git}{Git est un syst\`eme de contr\^ole de version
distribu\'e.}{L'ensemble du code (docker-compose, workers, configuration des
agents, rapport) est versionn\'e dans un d\'ep\^ot Git local, ce qui garantit
la tra\c cabilit\'e des modifications et la reproductibilit\'e du
d\'eploiement.}

\subsection{Tableau r\'ecapitulatif}

\begin{center}
\small
\begin{tabularx}{\linewidth}{@{}>{\bfseries}l l X@{}}
\toprule
\textbf{Couche} & \textbf{Technologie} & \textbf{R\^ole} \\
\midrule
Orchestration & Celery + Celery Beat & Pipeline de t\^aches asynchrones et ordonnancement \\
Courtier & RabbitMQ 3.12 & Files de messages distribu\'ees \\
\'Etat/cache & Redis 7 & \'Etat de scan, backend de r\'esultats, cache \\
Scanner & nmap (python-nmap) & D\'ecouverte + scan complet + OS \\
Intel menaces & MISP & Partage et enrichissement d'\'ev\'enements \\
Flux CVE & NVD API & Corr\'elation produit $\rightarrow$ CVE \\
Contexte exploit & EPSS, CISA KEV, OSV, Exploit-DB & Probabilit\'e d'exploitation et exploits \\
Moteur de risque & FastAPI & Calcul du score de risque \\
ITSM & GLPI & Tickets de rem\'ediation et gestion de parc \\
Recherche/viz & Elasticsearch + Kibana 8.11 & Indexation, tableaux de bord, alertes \\
Ingestion & Logstash 8.11 & Pipeline HTTP $\rightarrow$ Elasticsearch \\
Agents & Fleet Server / Elastic Agent / Beats & Collecte et gestion des agents \\
\bottomrule
\end{tabularx}
\end{center}
